%% kapa: exam=1 round=35 year=2024 period=2 form=A subject=appraisal_law subject_ko=감정평가관계법규 seq=11 number=11
%% source: Q-Net 감정평가사 기출문제 (cst00309) artlSeq=5213229 file=raw/1008/5213229/2교시 문제지 원본.pdf
%% conversion: kordoc 4.24.51 -> tools/gen1.py (auto); review=proofread-sonnet; flags=none
\begin{question}{11}
국토의 계획 및 이용에 관한 법령상 용도지역에 관한 설명으로 옳은 것은?
\begin{choices}
  \item 용도지역을 세분하는 지정은 도시ㆍ군기본계획으로도 할 수 있다.
  \item 하나의 시ㆍ도 안에서 둘 이상의 시ㆍ군에 걸쳐 지정되는 용도지역에 대해서는 국토교통부장관이 직접 도시ㆍ군관리계획을 입안할 수 있다.
  \item 하천의 매립목적이 그 매립구역과 이웃하고 있는 용도지역의 내용과 같으면 도시ㆍ군관리계획의 입안 및 결정 절차 없이 그 매립준공구역은 이웃하고 있는 용도지역으로 지정된 것으로 본다.
  \item 「산업입지 및 개발에 관한 법률」에 따라 국가산업단지로 지정된 지역은 「국토의 계획 및 이용에 관한 법률」에 따른 도시지역으로 결정ㆍ고시된 것으로 본다.
  \item 「택지개발촉진법」에 따른 택지개발지구가 개발사업의 완료로 해제되는 경우 그 지역은 택지개발지구를 지정하기 이전의 용도지역으로 환원된 것으로 본다.
\end{choices}
\end{question}
\answer{4}
